\documentclass[12pt]{article}
\usepackage{amsmath,amssymb,amsfonts}
\usepackage[margin=1in]{geometry}

\begin{document}

\textbf{Chapter 8, Problem 20.} Consider the quantity
\[
f^l(E) = \frac{2\tau}{(q_i^2+\mu^2)^l + \frac{\tau}{\pi}(q_i-i\mu)^3}
\]
where $\tau = 4\pi mg^2/\hbar^2$ and $q_i = (2mE_i/\hbar^2)^{1/2}$ [see Eq.\ (8.73)]. Show that if we consider $E$ to be a complex variable, that is, if we do not restrict $E$ to take on just physical real values, then $f^l(E)$ is a single-valued function of $E$ on a two-sheeted Riemann surface cut along the positive real axis. Show that, as a function of $E$, $f^l(E)$ is analytic everywhere except for a pole of second order at $E = -\hbar^2\mu^2/2m$ on the first sheet, a simple pole at $E = -\hbar^2\mu^2/2m$ on the second sheet (from the square root), and a so-called bound-state pole which may be either on the first or second sheet. If $\mu > \sqrt{\tau}$, then this pole is on the second sheet at $E = -\hbar^2(\sqrt{\tau}-\mu)^2/(2m)$; if $\sqrt{\tau}>\mu$, the pole is on the first sheet at $E = -\hbar^2(\sqrt{\tau}-\mu)^2/(2m)$. This value of $E$ is precisely the energy of the bound state determined in the previous problem.

\bigskip
\textbf{Solution.}

The partial-wave scattering amplitude for the separable potential $V = -g^2|\phi\rangle\langle\phi|$ with $\phi(r) = e^{-\mu r}$ is given by (from the $T$-matrix formalism):
\[
f^l(E) = -\frac{2m}{\hbar^2}\frac{g^2|\tilde\phi(q)|^2}{1-g^2\Pi(E)},
\]
where $\Pi(E) = \frac{2m}{\hbar^2}\int\frac{|\tilde\phi(k)|^2}{k^2-q^2-i\epsilon}\frac{k^2\,dk}{2\pi^2}$ and $q = \sqrt{2mE/\hbar^2}$.

\textbf{Analytic structure:}

The key variable is $q = \sqrt{2mE/\hbar^2}$. For real positive $E$, $q$ is real and positive. But as a function of complex $E$, $q$ has a branch point at $E=0$, requiring a branch cut (conventionally along the positive real axis) and creating a two-sheeted Riemann surface:
\begin{itemize}
\item \textbf{First (physical) sheet:} $\text{Im}\,q > 0$ (for $E$ in the upper half-plane).
\item \textbf{Second (unphysical) sheet:} $\text{Im}\,q < 0$.
\end{itemize}

The function $f^l(E)$ inherits this two-sheeted structure through its dependence on $q$.

\textbf{Poles:}

1. The form factor $\tilde\phi(q) \propto 1/(q^2+\mu^2)^2$ has poles at $q^2 = -\mu^2$, i.e., $q = \pm i\mu$. On the first sheet ($\text{Im}\,q>0$), this corresponds to $q = i\mu$, giving $E = -\hbar^2\mu^2/(2m)$. This is a \textbf{pole of order 2} in $f^l$ (from the denominator of $|\tilde\phi|^2$, which involves $(q^2+\mu^2)^4$---but only those factors not cancelled contribute). Actually, the pole comes from the form factor squared, giving a second-order pole at $E = -\hbar^2\mu^2/(2m)$ on the first sheet.

2. On the second sheet, $q = -i\mu$ gives the same energy $E = -\hbar^2\mu^2/(2m)$ but as a simple pole.

3. The \textbf{bound-state pole} comes from the zero of the denominator $1-g^2\Pi(E) = 0$. After evaluation, this occurs at:
\[
E_{\text{pole}} = -\frac{\hbar^2}{2m}(\sqrt{\tau}-\mu)^2,
\]
where $\tau = 4\pi mg^2/\hbar^2$.

\begin{itemize}
\item If $\sqrt{\tau}>\mu$: The pole is at $q = i(\sqrt{\tau}-\mu)$ with $\text{Im}\,q > 0$, so it is on the \textbf{first sheet}. This represents a genuine bound state.
\item If $\mu>\sqrt{\tau}$: The pole is at $q = -i(\mu-\sqrt{\tau})$ with $\text{Im}\,q < 0$, so it is on the \textbf{second sheet}. This is a virtual state (not a true bound state).
\end{itemize}

The condition for a bound state ($\sqrt{\tau}>\mu$) is equivalent to $g^2 > \hbar^2\mu^2/(4\pi m)$, which matches Problem 19. The bound-state energy $E_B = \hbar^2(\sqrt{\tau}-\mu)^2/(2m)$ also matches.

As $g^2$ drops below the critical value, the bound-state pole moves through the threshold $E=0$ from the first sheet to the second sheet, becoming a virtual state. The effect of the bound state does not magically vanish from the scattering amplitude; rather, it moves off the physical sheet. $\blacksquare$

\end{document}
